\documentclass[../../main.tex]{subfiles}
\begin{document}
\section{淹没}\label{section:ch04-淹没}

秩定理最重要的应用之一,是极大地拓展我们对淹没性质的理解。若 $\pi:M\to N$ 是任一连续映射, $\pi$ 的一个\textbf{截面}是 $\pi$ 的连续右逆,即连续映射 $\sigma:N\to M$ 满足
\begin{align*}
\pi\circ\sigma=\operatorname{Id}_N.
\end{align*}
$\pi$ 的一个\textbf{局部截面}是定义在某开子集 $U\subseteq N$ 上、并满足类似关系 $\pi\circ\sigma=\operatorname{Id}_U$ 的连续映射 $\sigma:U\to M$。光滑淹没许多重要性质都源于这样一个事实:它们容许丰富的光滑局部截面。

\begin{theorem}[Local Section Theorem]\label{theorem:ch04-local-section}
设 $M$ 与 $N$ 为光滑流形, $\pi:M\to N$ 为光滑映射。则 $\pi$ 是光滑淹没,当且仅当 $M$ 的每一点都落在 $\pi$ 的某个光滑局部截面的像中。
\end{theorem}

\begin{proof}
先设 $\pi$ 是光滑淹没。给定 $p\in M$,令 $q=\pi(p)\in N$。由秩定理,我们可选取以 $p$ 为中心的坐标 $\bigl(x^1,\dots,x^m\bigr)$ 与以 $q$ 为中心的坐标 $\bigl(y^1,\dots,y^n\bigr)$,使 $\pi$ 具有坐标表示
\begin{align*}
\pi\bigl(x^1,\dots,x^n,x^{n+1},\dots,x^m\bigr)=\bigl(x^1,\dots,x^n\bigr).
\end{align*}
若 $\varepsilon$ 是充分小的正数,则坐标立方体
\begin{align*}
C_\varepsilon=\{x:|x^i|<\varepsilon,\ i=1,\dots,m\}
\end{align*}
是 $p$ 的邻域,其在 $\pi$ 下的像是立方体
\begin{align*}
C'_\varepsilon=\{y:|y^i|<\varepsilon,\ i=1,\dots,n\}.
\end{align*}
坐标表示为
\begin{align*}
\widehat{\sigma}\bigl(x^1,\dots,x^n\bigr)=\bigl(x^1,\dots,x^n,0,\dots,0\bigr)
\end{align*}
的映射 $\sigma:C'_\varepsilon\to C_\varepsilon$ 是 $\pi$ 的光滑局部截面,满足 $\sigma(q)=p$(\cref{figure:ch04-local-section})。

反之,设 $M$ 的每一点都落在某光滑局部截面的像中。给定 $p\in M$,令 $\sigma:U\to M$ 为光滑局部截面使 $\sigma(q)=p$,其中 $q=\pi(\sigma(q))=\pi(p)\in N$。等式 $\pi\circ\sigma=\operatorname{Id}_U$ 蕴涵 $\mathrm{d}\pi_p\circ\mathrm{d}\sigma_q=\operatorname{Id}_{T_qN}$,而这又蕴涵 $\mathrm{d}\pi_p$ 满射。
\end{proof}

\begin{figure}[ht]
\centering
\includegraphics[width=0.8\textwidth]{fig-04-4.png}
\caption{淹没的局部截面}\label{figure:ch04-local-section}
\end{figure}

这个定理引出如下定义:若 $\pi:X\to Y$ 是连续映射,称 $\pi$ 是\textbf{拓扑淹没},如果 $X$ 的每一点都落在 $\pi$ 的某(连续)局部截面的像中。上一定理表明每个光滑淹没都是拓扑淹没。

\begin{exercise}\label{exercise:ch04-topological-not-smooth}
举一个光滑映射的例子,它是拓扑淹没但不是光滑淹没。
\end{exercise}

\begin{proposition}[Properties of Smooth Submersions]\label{proposition:ch04-submersion-open-quotient}
设 $M$ 与 $N$ 为光滑流形,并设 $\pi:M\to N$ 是光滑淹没。则 $\pi$ 是开映射;若它满射,则它是商映射。
\end{proposition}

\begin{proof}
设 $W$ 是 $M$ 的开子集, $q$ 是 $\pi(W)$ 中的一点。对任一满足 $\pi(p)=q$ 的 $p\in W$,存在 $q$ 的邻域 $U$,使 $\pi$ 在其上有光滑局部截面 $\sigma:U\to M$ 满足 $\sigma(q)=p$。对每个 $y\in\sigma^{-1}(W)$, $\sigma(y)\in W$ 这一事实蕴涵 $y=\pi(\sigma(y))\in\pi(W)$。于是 $\sigma^{-1}(W)$ 是包含在 $\pi(W)$ 中的 $q$ 的邻域,这蕴涵 $\pi(W)$ 开。第二个断言由第一个推出,因为每个满射开连续映射都是商映射。
\end{proof}

下面三个定理提供了我们研究淹没时将频繁使用的重要工具。请注意这些结果与\cref{theorem:appA-quotient-char}、\cref{theorem:appA-passing-quotient} 以及\cref{theorem:appA-uniqueness-quotient} 的相似之处。这表明满射光滑淹没在光滑流形理论中所起的作用,类似于商映射在拓扑学中所起的作用。第一个定理推广了习题 4.10(b) 的结果。

\begin{theorem}[Characteristic Property of Surjective Smooth Submersions]\label{theorem:ch04-char-submersion}
设 $M$ 与 $N$ 为光滑流形, $\pi:M\to N$ 是满射光滑淹没。对任一(带边或不带边的)光滑流形 $P$,映射 $F:N\to P$ 光滑当且仅当 $F\circ\pi$ 光滑:
\begin{align*}
\begin{array}{ccc}
M & \xrightarrow{\ F\ } & P\\
\pi\bigg\downarrow &&\\
N
\end{array}
\end{align*}
\end{theorem}

\begin{proof}
若 $F$ 光滑,则由复合知 $F\circ\pi$ 光滑。反之,设 $F\circ\pi$ 光滑,并令 $q\in N$ 任意。由于 $\pi$ 满射,存在点 $p\in\pi^{-1}(q)$,于是局部截面定理保证存在 $q$ 的邻域 $U$ 与 $\pi$ 的光滑局部截面 $\sigma:U\to M$ 使 $\sigma(q)=p$。于是 $\pi\circ\sigma=\operatorname{Id}_U$ 蕴涵
\begin{align*}
F|_U=F|_U\circ\operatorname{Id}_U=F|_U\circ(\pi\circ\sigma)=(F\circ\pi)\circ\sigma,
\end{align*}
这是光滑映射的复合。这表明 $F$ 在每一点附近都光滑,故它光滑。
\end{proof}

习题 4-7 说明了这一性质在何种意义下是``特征性的''。下一定理给出一个非常一般的充分条件,使一个光滑映射能被淹没``推下去''。

\begin{theorem}[Passing Smoothly to the Quotient]\label{theorem:ch04-pass-quotient}
设 $M$ 与 $N$ 为光滑流形, $\pi:M\to N$ 是满射光滑淹没。若 $P$ 是带边或不带边的光滑流形, $F:M\to P$ 是在 $\pi$ 的各纤维上取常值的光滑映射,则存在唯一光滑映射 $\widetilde{F}:N\to P$ 使 $\widetilde{F}\circ\pi=F$:
\begin{align*}
\begin{array}{ccc}
M & \xrightarrow{\ F\ } & P\\
\pi\bigg\downarrow & \nearrow &\\
N & \widetilde{F} &
\end{array}
\end{align*}
\end{theorem}

\begin{proof}
由于满射光滑淹没是商映射, \cref{theorem:appA-passing-quotient} 表明存在唯一连续映射 $\widetilde{F}:N\to P$ 满足 $\widetilde{F}\circ\pi=F$。它由\cref{theorem:ch04-char-submersion} 知光滑。
\end{proof}

最后,我们有如下唯一性结果。

\begin{theorem}[Uniqueness of Smooth Quotients]\label{theorem:ch04-uniqueness-quotient}
设 $M,N_1,N_2$ 为光滑流形, $\pi_1:M\to N_1$ 与 $\pi_2:M\to N_2$ 是在彼此的纤维上取常值的满射光滑淹没。则存在唯一微分同胚 $F:N_1\to N_2$ 使 $F\circ\pi_1=\pi_2$:
\begin{align*}
\begin{array}{ccc}
 & M &\\
\pi_1\swarrow & \searrow\pi_2 &\\
N_1 & \xrightarrow{\ F\ } & N_2
\end{array}
\end{align*}
\end{theorem}

\begin{exercise}\label{exercise:ch04-prove-uniqueness-quotient}
证明\cref{theorem:ch04-uniqueness-quotient}。
\end{exercise}

\end{document}